
\documentclass[11pt,a4paper]{article}

\usepackage[T1]{fontenc}
\usepackage[utf8]{inputenc}
\usepackage{lmodern}
\usepackage[margin=1in]{geometry}
\usepackage{amsmath,amssymb,amsthm,mathtools}
\usepackage{booktabs}
\usepackage{array}
\usepackage{tabularx}
\usepackage{microtype}
\usepackage{enumitem}
\usepackage[hidelinks]{hyperref}
\usepackage[nameinlink,noabbrev]{cleveref}

\setlength{\parskip}{0.45em}
\setlength{\parindent}{0pt}

\newtheorem{theorem}{Theorem}[section]
\newtheorem{corollary}[theorem]{Corollary}
\newtheorem{lemma}[theorem]{Lemma}
\newtheorem{proposition}[theorem]{Proposition}
\theoremstyle{definition}
\newtheorem{definition}[theorem]{Definition}
\newtheorem{assumption}[theorem]{Assumption}
\theoremstyle{remark}
\newtheorem*{remark}{Remark}

\DeclareMathOperator{\Spec}{Spec}
\DeclareMathOperator{\Fix}{Fix}
\DeclareMathOperator{\tr}{Tr}
\DeclareMathOperator{\SF}{SF}

\newcolumntype{Y}{>{\raggedright\arraybackslash}X}

\newcommand{\Seam}{\Sigma}
\newcommand{\inv}{\iota}
\newcommand{\Tmin}{\mathfrak T_{\mathrm{min}}}
\newcommand{\Tprim}{\mathfrak T_{\mathrm{prim}}}
\newcommand{\cthree}{c_3}
\newcommand{\bone}{b_1}
\newcommand{\phibase}{\varphi_{\mathrm{base}}}
\newcommand{\phiseam}{\varphi_{\Sigma}}
\newcommand{\phiz}{\varphi_{0}^{\mathrm{ret}}}
\newcommand{\deltatop}{\delta_{\mathrm{top}}}
\newcommand{\gammaY}{\gamma}
\newcommand{\SM}{\mathrm{SM}}
\newcommand{\Oadm}{\Omega_{\mathrm{adm}}}

\numberwithin{equation}{section}

\title{\textbf{TFPT Minimal Kernel}\\\large Clean core note for downstream application papers}
\author{Stefan Hamann \and Alessandro Rizzo}
\date{\today}

\begin{document}
\maketitle

\begin{abstract}
This note isolates the smallest TFPT theorem packet needed as a reusable core.
The retained layer consists of the admissible reduced datum, the reconstructed primitive scaffold,
the unique minimal carrier $E_3 \oplus E_2$, the hypercharge--seam inversion,
the global Standard Model quotient
\[
S(U(3)\times U(2)) \cong \frac{SU(3)\times SU(2)\times U(1)_Y}{\mathbb Z_6},
\]
one chiral family with $\nu^c$ from the positive half spinor, and the exact electromagnetic
closure equation on the canonical branch. Everything else is intentionally left outside this
kernel so that short downstream notes can extend a stable core without dragging the full theorem
surface behind them.
\end{abstract}

\tableofcontents

\section{Reduced datum and primitive scaffold}

\begin{definition}[Admissible reduced theorem datum]
\label{def:reduced-datum}
The admissible reduced datum is
\[
\Tmin^{\mathrm{adm}} := (\mathcal A,\mathcal H,D,J,\Gamma,\inv,\nu)
\]
with the following built in conditions:
\begin{enumerate}[label=(\roman*), leftmargin=1.6em, itemsep=0.25em]
    \item $Z(\mathcal A)$ is spectrally regular and $\Spec Z(\mathcal A)$ is a smooth compact oriented manifold $\widetilde M$;
    \item $\inv$ acts as a smooth involution of order two on $Z(\mathcal A)$, commutes with $J$ and $\Gamma$, and has fixed locus
    \[
    \Seam := \Fix(\inv);
    \]
    \item in a collar neighborhood of $\Seam$, the Dirac operator has product form
    \[
    D = \gamma_n(\partial_n + B_{\Seam});
    \]
    \item the cutoff profile is generated by a positive complete Bernstein measure
    \[
    f(x) = \int_0^\infty e^{-t x^2}\,d\nu(t),
    \qquad
    d\nu(t)\ge 0;
    \]
    \item the seam collar is reflection regular, so the doubled operator and the later admissibility complex are well posed on the same datum.
\end{enumerate}
\end{definition}

\begin{assumption}[Almost commutative factorization]
\label{ass:factorization}
The reduced datum admits a factorization
\[
\mathcal H \cong L^2(\widetilde M,S)\otimes \mathcal H_{\mathrm{int}},
\]
where $S\to \widetilde M$ is the spinor bundle on the oriented double cover and
$\mathcal H_{\mathrm{int}}$ is a finite internal Hilbert space.
\end{assumption}

\begin{theorem}[Primitive scaffold reconstructed from the reduced datum]
\label{thm:primitive-scaffold}
Let $\Tmin^{\mathrm{adm}}$ be as in \Cref{def:reduced-datum}. Then the primitive kinematic scaffold is reconstructed canonically by
\[
\widetilde M = \Spec Z(\mathcal A),
\qquad
M = \widetilde M/\inv,
\qquad
\Seam = \Fix(\inv),
\qquad
X = \widetilde M \setminus \Seam,
\]
together with the reference and relative operators
\[
D_{\mathrm{ref}} := \frac12(D+\inv D\inv),
\qquad
D_{\mathrm{rel}} := \frac12(D-\inv D\inv).
\]
Hence
\[
\Tprim^{\mathrm{kin}}
:=
(X,\widetilde M,M,\Seam,\inv,D_{\mathrm{ref}},D_{\mathrm{rel}},f)
\]
is derived rather than assumed. Under \Cref{ass:factorization}, the spinor bundle and the internal
Hilbert space refer to the chosen factorization and are not additional independent inputs.
\end{theorem}

\begin{remark}
This note deliberately stops at the reconstructed primitive scaffold. The admissibility selector,
geometric completion, transport sector, and observation layer belong to later TFPT layers and are
left out here on purpose.
\end{remark}

\section{Minimal carrier theorem}

\begin{theorem}[Minimal carrier theorem]
\label{thm:minimal-carrier}
Let $E = E_b \oplus E_s$ be a split carrier with traceless diagonal generator
\[
Y_{b,s}
=
\mathrm{diag}\!\left(
\underbrace{-\frac1b,\ldots,-\frac1b}_{b\text{ times}},
\underbrace{\frac1s,\ldots,\frac1s}_{s\text{ times}}
\right).
\]
If one simultaneously requires
\[
\dim \Lambda^{\mathrm{even}} E = 16,
\qquad
\tr_E Y_{b,s}^2 = \frac56,
\]
then necessarily $\{b,s\} = \{3,2\}$.
\end{theorem}

\begin{proof}
Set $g := b+s$. Since $\dim \Lambda^{\mathrm{even}}E = 2^{g-1}$, the first condition gives $g=5$.
For the displayed diagonal generator,
\[
\tr_E Y_{b,s}^2
=
b\left(\frac1{b^2}\right)+s\left(\frac1{s^2}\right)
=
\frac1b+\frac1s
=
\frac{b+s}{bs}
=
\frac{5}{bs}.
\]
Imposing $\tr_E Y_{b,s}^2 = 5/6$ yields $bs=6$. Together with $b+s=5$, this forces the unordered
solution $\{b,s\}=\{3,2\}$.
\end{proof}

From now on we therefore fix
\[
E = E_3 \oplus E_2,
\qquad
P_3 = \frac{\mathbf 1-\inv}{2},
\qquad
P_2 = \frac{\mathbf 1+\inv}{2},
\]
and read hypercharge as
\begin{equation}
\label{eq:hypercharge}
Y
=
\mathrm{diag}\!\left(-\frac13,-\frac13,-\frac13,\frac12,\frac12\right)
=
-\frac13 P_3+\frac12 P_2
=
\frac1{12}\mathbf 1+\frac5{12}\inv.
\end{equation}

\begin{lemma}[Hypercharge--seam inversion]
\label{lem:hypercharge-seam}
On the minimal carrier,
\begin{equation}
\label{eq:seam-inversion}
\inv = \frac{12Y-\mathbf 1}{5},
\end{equation}
and therefore
\begin{equation}
\label{eq:carrier-polynomial}
6Y^2 - Y - \mathbf 1 = 0.
\end{equation}
Consequently,
\[
P_2 = \frac{\mathbf 1+\inv}{2} = \frac{2(3Y+\mathbf 1)}{5},
\qquad
P_3 = \frac{\mathbf 1-\inv}{2} = \frac{3(\mathbf 1-2Y)}{5}.
\]
\end{lemma}

\begin{proof}
Invert the affine relation in \Cref{eq:hypercharge}. Squaring \Cref{eq:seam-inversion} and using
$\inv^2=\mathbf 1$ gives \Cref{eq:carrier-polynomial}. The projector formulas follow from
$P_2=(\mathbf 1+\inv)/2$ and $P_3=(\mathbf 1-\inv)/2$.
\end{proof}

\begin{corollary}[Carrier invariants]
\label{cor:carrier-invariants}
On the minimal branch,
\[
\gammaY := \tr_E Y^2 = \frac56,
\qquad
B := \frac{\mathrm{rank}\,E_3}{\mathrm{rank}\,E_2} = \frac32,
\qquad
B\gammaY = \frac54.
\]
\end{corollary}

\begin{remark}
Because of \Cref{eq:carrier-polynomial}, every analytic function of $Y$ collapses to an affine
combination of $\mathbf 1$ and $Y$. In that precise sense, the hard carrier is a two point algebra
supported on the spectral values $-1/3$ and $1/2$.
\end{remark}

\section{Global quotient and one family}

\begin{theorem}[Global Standard Model theorem]
\label{thm:global-sm}
The determinant preserving connected stabilizer of the minimal carrier split is
\[
G_{\SM}^{\mathrm{glob}}
=
S(U(3)\times U(2))
\cong
\frac{SU(3)\times SU(2)\times U(1)_Y}{\mathbb Z_6}.
\]
\end{theorem}

\begin{proof}
Any automorphism preserving the split $E_3\oplus E_2$ acts by a pair $(u_3,u_2)\in U(3)\times U(2)$.
Determinant preservation on the full carrier imposes $\det(u_3)\det(u_2)=1$, hence the connected
stabilizer is exactly $S(U(3)\times U(2))$.

Now define
\[
\Phi : SU(3)\times SU(2)\times U(1)_Y \longrightarrow S(U(3)\times U(2)),
\qquad
\Phi(g_3,g_2,z) := (z^2 g_3, z^{-3} g_2).
\]
The determinant condition is automatic since $z^6 z^{-6}=1$. Surjectivity follows by choosing
$z\in U(1)$ with $z^6=\det(u_3)=\det(u_2)^{-1}$ and setting
\[
g_3 := z^{-2}u_3 \in SU(3),
\qquad
g_2 := z^3 u_2 \in SU(2).
\]
The kernel consists of
\[
\left\{
\bigl(z^{-2}\mathbf 1_3,\ z^3\mathbf 1_2,\ z\bigr): z^6=1
\right\}
\cong \mathbb Z_6,
\]
which proves the quotient description.
\end{proof}

Matter is read from the positive half spinor
\[
S^+ := \Lambda^{\mathrm{even}}E,
\qquad
\dim S^+ = 16.
\]

\begin{theorem}[One carrier family]
\label{thm:one-family}
The positive half spinor of the minimal carrier decomposes as
\[
S^+
\cong
(1,1)_0\oplus(3,2)_{1/6}\oplus(\bar 3,1)_{-2/3}\oplus(\bar 3,1)_{1/3}\oplus(1,2)_{-1/2}\oplus(1,1)_1,
\]
that is, exactly one chiral Standard Model family with $\nu^c$. Equivalently,
\[
S^+ \cong \mathbf 1 \oplus \mathbf{10} \oplus \overline{\mathbf 5}.
\]
\end{theorem}

\begin{proof}
Since
\[
S^+ = \Lambda^0E \oplus \Lambda^2E \oplus \Lambda^4E,
\]
the sectors are read degree by degree. One has
\[
\Lambda^0E \cong (1,1)_0.
\]
For the minimal carrier,
\[
\Lambda^2E
\cong
\Lambda^2E_3 \oplus (E_3\otimes E_2)\oplus \Lambda^2E_2
\cong
(\bar 3,1)_{-2/3}\oplus(3,2)_{1/6}\oplus(1,1)_1.
\]
Because the full determinant is trivial in $S(U(3)\times U(2))$, one has $\Lambda^5E\cong 1$ and
hence $\Lambda^4E\cong E^\ast$, namely
\[
\Lambda^4E \cong (1,2)_{-1/2}\oplus(\bar 3,1)_{1/3}.
\]
Collecting the three even degree sectors yields the stated one family packet with total
multiplicity $1+6+3+3+2+1=16$.
\end{proof}

\begin{table}[ht]
\centering
\small
\begin{tabularx}{\linewidth}{@{}>{\raggedright\arraybackslash}p{0.20\linewidth}>{\raggedright\arraybackslash}p{0.20\linewidth}>{\centering\arraybackslash}p{0.12\linewidth}>{\raggedright\arraybackslash}X>{\centering\arraybackslash}p{0.10\linewidth}@{}}
\toprule
Exterior component & $\SM$ representation & Hypercharge & Field content & Mult. \\
\midrule
$\Lambda^0E$ & $(1,1)_0$ & $0$ & $\nu^c$ & $1$ \\
$\Lambda^2E_3$ & $(\bar 3,1)_{-2/3}$ & $-2/3$ & $u^c$ & $3$ \\
$E_3\otimes E_2$ & $(3,2)_{1/6}$ & $1/6$ & $Q_L$ & $6$ \\
$\Lambda^2E_2$ & $(1,1)_1$ & $1$ & $e^c$ & $1$ \\
$\Lambda^3E_3\otimes E_2$ & $(1,2)_{-1/2}$ & $-1/2$ & $L_L$ & $2$ \\
$\Lambda^2E_3\otimes\Lambda^2E_2$ & $(\bar 3,1)_{1/3}$ & $1/3$ & $d^c$ & $3$ \\
\midrule
\multicolumn{4}{r}{Total multiplicity} & $16$ \\
\bottomrule
\end{tabularx}
\caption{One chiral family from $S^+=\Lambda^{\mathrm{even}}E$, written in left chiral notation with charge conjugate singlets.}
\label{tab:one-family}
\end{table}

\section{Electromagnetic closure and kernel package}

\begin{proposition}[Seam normalization and base opening]
\label{prop:seam-normalization}
If the seam loop carries one unit of relative spectral flow,
\[
\SF(U_{\Seam}) = 1,
\qquad
\Delta_{\Seam} = 2\pi,
\]
then the seam constant is
\[
\cthree := \frac{\Delta_{\Seam}}{16\pi^2} = \frac{1}{8\pi}.
\]
Moreover the seam even base opening is
\[
\phibase
:=
\frac{1-\gammaY}{\pi}
=
\frac{1-\tr_EY^2}{\pi}
=
\frac{1}{6\pi}.
\]
\end{proposition}

\begin{definition}[Canonical abelian coefficient]
\label{def:b1}
For family number $N_{\mathrm{fam}}$ and light Higgs doublet count $N_\Phi$, define
\[
\bone(N_{\mathrm{fam}},N_\Phi) := \frac43 N_{\mathrm{fam}} + \frac1{10}N_\Phi.
\]
On the canonical branch,
\[
N_{\mathrm{fam}}=3,
\qquad
N_\Phi=1,
\qquad
\bone = \frac{41}{10}.
\]
\end{definition}

\begin{definition}[Exact seam generating function]
\label{def:seam-generating-function}
Set
\[
\deltatop := \frac{3}{256\pi^4},
\qquad
x := e^{-2\alpha}.
\]
The exact seam opening is
\begin{equation}
\label{eq:phi-seam}
\phiseam(\alpha)
:=
\phibase + \deltatop x(1-\deltatop x)^{-5/4}.
\end{equation}
Its expansion on the admissible branch begins as
\[
\phiseam(\alpha)
=
\phibase
+\deltatop e^{-2\alpha}
+\frac54 \deltatop^2 e^{-4\alpha}
+O(e^{-6\alpha}).
\]
\end{definition}

\begin{theorem}[Exact electromagnetic closure]
\label{thm:electromagnetic-closure}
On the canonical admissible branch, the fine structure constant is the unique positive solution of
\begin{equation}
\label{eq:carrier-cfe}
\alpha^3 - 2\cthree^3\alpha^2 - 8\cthree^6\,\bone(N_{\mathrm{fam}},N_\Phi)\,
\ln\!\bigl((\phiseam(\alpha))^{-1}\bigr) = 0.
\end{equation}
For the canonical family branch $N_{\mathrm{fam}}=3$ and $N_\Phi=1$, this becomes
\begin{equation}
\label{eq:carrier-cfe-canonical}
\alpha^3 - \frac{\alpha^2}{256\pi^3}
-
\frac{41}{327680\pi^6}\ln\!\bigl((\phiseam(\alpha))^{-1}\bigr)
= 0.
\end{equation}
\end{theorem}

\begin{corollary}[Kernel output package]
\label{cor:kernel-package}
The minimal electromagnetic package consists of
\[
\cthree = \frac{1}{8\pi},
\qquad
\phiz = \frac{1}{6\pi}+\frac{3}{256\pi^4},
\qquad
\alpha^{-1}(0)=137.03599921615843,
\]
together with the seed chain
\[
C_{a\gamma\gamma} = -4\cthree = -\frac{1}{2\pi},
\qquad
\beta_{\mathrm{rad}} = \frac{\phiz}{4\pi},
\qquad
\Omega_b = (4\pi-1)\beta_{\mathrm{rad}},
\]
\[
\lambda_C = \sqrt{\phiz(1-\phiz)},
\qquad
\sin^2\theta_{13} = \phiz e^{-\gammaY}.
\]
\end{corollary}

\begin{remark}
The $\alpha$ line must use the exact self consistent opening $\phiseam(\alpha)$ from
\Cref{eq:phi-seam}. The retained seed $\phiz$ is only the frozen bridge value used for the
one seed relations in \Cref{cor:kernel-package}.
\end{remark}

\appendix

\section{Optional arithmetic compression on the canonical branch}

The following identity is not needed for the minimal carrier theorem itself, but it is a compact
bridge relation that is often useful in short downstream notes. Define the occupancy count
\[
\Omega_{\mathrm{occ}} := 16\,N_{\mathrm{fam}}.
\]
Then
\[
10\,\bone
=
\Omega_{\mathrm{occ}}\,\gammaY + N_\Phi.
\]
On the canonical branch,
\[
\Omega_{\mathrm{occ}} = \Oadm = 48,
\qquad
\gammaY = \frac56,
\qquad
N_\Phi = 1,
\]
so the hard chain becomes
\[
10\,\bone = \Oadm\gammaY + 1 = 41.
\]

\section{Minimal import block for application notes}

For downstream sector papers, the imported TFPT kernel can be summarized in one line as
\[
\bigl(
E_3\oplus E_2,\,
Y,\,
\inv,\,
G_{\SM}^{\mathrm{glob}},\,
S^+,\,
\gammaY,\,
B,\,
\cthree,\,
\phibase,\,
\deltatop,\,
\bone,\,
\phiseam(\alpha)
\bigr).
\]
That block is intentionally small enough to be copied into a short application note without
pulling in the full closure architecture.

\end{document}
